\chapter{От текста до работающей схемы}

\section{Маршрут проектирования}

Для программиста путь от программы до работающей прошивки выглядит так:
«исходный код $\to$ компилятор $\to$ машинные команды». У разработчика ПЛИС
шагов больше, и каждый решает свою задачу.

\begin{figure}[H]
\centering
\begin{tikzpicture}[node distance=4mm and 11mm,
  fstep/.style={block, minimum width=26mm, minimum height=15mm},
  file/.style={draw=FpgaGray, fill=white, font=\ttfamily\scriptsize, minimum height=5mm, inner sep=2pt},
]
  \node[fstep] (hdl) {\textbf{1. Описание}\\\scriptsize Verilog (.v)};
  \node[fstep, right=of hdl] (syn) {\textbf{2. Синтез}\\\scriptsize Yosys / GowinSynth};
  \node[fstep, right=of syn] (pnr) {\textbf{3. Размещение}\\\textbf{и трассировка}\\\scriptsize nextpnr / Gowin PnR};
  \node[fstep, right=of pnr] (bit) {\textbf{4. Битстрим}\\\scriptsize gowin\_pack};
  \node[hblock, minimum width=26mm, minimum height=15mm, below=12mm of bit] (load) {\textbf{5. Загрузка}\\\scriptsize openFPGALoader\\\scriptsize Gowin Programmer};
  \node[oblock, minimum width=26mm, minimum height=15mm, left=of load] (board) {\textbf{Tang Nano 20K}};
  \node[gblock, minimum width=26mm, minimum height=15mm, below=12mm of hdl] (sim) {\textbf{Симуляция}\\\scriptsize Icarus Verilog\\\scriptsize GTKWave};
  \node[file, above=7mm of pnr] (cst) {pins.cst};
  \draw[arr] (hdl) -- (syn);
  \draw[arr] (syn) -- node[below, lbl] {netlist} (pnr);
  \draw[arr] (pnr) -- (bit);
  \draw[arr] (bit) -- node[right, lbl] {.fs} (load);
  \draw[arr] (load) -- node[above, lbl] {USB} (board);
  \draw[arr, <->, dashed] (hdl) -- node[right, lbl] {проверка} (sim);
  \draw[arr] (cst) -- (pnr);
  \node[lbl, right=1mm of cst] {номера ножек};
\end{tikzpicture}
\caption{Маршрут проектирования для ПЛИС}
\label{fig:flow}
\end{figure}

\begin{enumerate}
  \item \textbf{Описание.} Вы пишете, \emph{какой должна быть схема}, на языке
        описания аппаратуры (HDL, Hardware Description Language). Мы будем
        использовать самый распространённый из них --- \textbf{Verilog}.
  \item \textbf{Синтез.} Программа-синтезатор читает текст и превращает его в
        \term{список соединений} (netlist): «вот такие LUT, вот такие триггеры,
        вот так соединены». При этом логика упрощается, как при минимизации
        по картам Карно, только автоматически.
  \item \textbf{Размещение и трассировка} (Place \& Route). Каждый LUT и
        триггер из списка получает конкретное место на кристалле, а соединения
        прокладываются по трассировочным каналам. Здесь же используется файл
        \code{.cst}, чтобы порты схемы попали на нужные ножки.
  \item \textbf{Битстрим.} Результат упаковывается в файл конфигурации
        \code{.fs}: те самые биты для каждого LUT и каждого ключа.
  \item \textbf{Загрузка.} Файл по USB передаётся в ПЛИС (в SRAM) или во
        Flash-память на плате.
\end{enumerate}

\section{Что делает синтезатор}

Рассмотрим, как синтезатор превращает текст в «железо». Пусть мы написали:
\begin{vcode}
assign y = (a & b) | ~c;
\end{vcode}

Человек нарисовал бы схему из трёх элементов: И, НЕ и ИЛИ. Но у ПЛИС нет
отдельных элементов --- есть LUT. Функция зависит от трёх переменных,
поэтому синтезатор вычисляет её таблицу истинности и записывает её в
\emph{один} LUT3 (или LUT4 с одним неиспользуемым входом).

\begin{figure}[H]
\centering
\begin{tikzpicture}
  \begin{scope}
    \node[gAND] (and) at (1.3,0.9) {};
    \node[gNOT] (not) at (1.1,-0.5) {};
    \node[gOR] (or) at (3.0,0.2) {};
    \draw[wire] (and.input 1) -- ++(-0.7,0) node[left] {$a$};
    \draw[wire] (and.input 2) -- ++(-0.7,0) node[left] {$b$};
    \draw[wire] (not.input) -- ++(-0.5,0) node[left] {$c$};
    \draw[wire] (and.output) -- ++(0.3,0) |- (or.input 1);
    \draw[wire] (not.output) -- ++(0.5,0) |- (or.input 2);
    \draw[wire] (or.output) -- ++(0.5,0) node[right] {$y$};
    \node[lbl] at (1.8,-1.5) {как задумал человек};
  \end{scope}
  \draw[-{Stealth[length=4mm]}, line width=2pt, FpgaOrange] (4.9,0.2) -- (6.1,0.2) node[midway, above=1mm, lbl] {синтез};
  \begin{scope}[xshift=7cm]
    \node[draw=FpgaBlue, fill=FpgaLight, thick, minimum width=18mm, minimum height=20mm, font=\sffamily\small, align=center] (lut) at (1.5,0.2) {LUT3};
    \draw[wire] ([yshift=5mm]lut.west) -- ++(-0.7,0) node[left] {$a$};
    \draw[wire] (lut.west) -- ++(-0.7,0) node[left] {$b$};
    \draw[wire] ([yshift=-5mm]lut.west) -- ++(-0.7,0) node[left] {$c$};
    \draw[wire] (lut.east) -- ++(0.7,0) node[right] {$y$};
    \node[lbl] at (1.5,-1.5) {как это реализует ПЛИС};
  \end{scope}
  \begin{scope}[xshift=11.2cm, yshift=1.6cm]
    \matrix[matrix of nodes, nodes={font=\ttfamily\scriptsize, inner sep=1.2pt},
            row 1/.style={nodes={font=\sffamily\bfseries\scriptsize}}] (tt) {
      a & b & c & y \\ 0&0&0&1 \\ 0&0&1&0 \\ 0&1&0&1 \\ 0&1&1&0 \\ 1&0&0&1 \\ 1&0&1&0 \\ 1&1&0&1 \\ 1&1&1&1 \\};
    \draw[FpgaGray] (tt-1-1.south west) -- (tt-1-4.south east);
    \node[lbl, below=0mm of tt] {память LUT};
  \end{scope}
\end{tikzpicture}
\caption{Синтез: три логических элемента превращаются в один LUT3 с
таблицей истинности внутри}
\end{figure}

\begin{ideabox}
Синтезатору всё равно, как вы нарисовали бы схему на бумаге. Важно лишь,
\emph{какую функцию} вы описали. Поэтому на Verilog описывают
\emph{поведение} схемы, а подбором элементов занимается программа.
\end{ideabox}

\section{Инструменты}

Для Tang Nano 20K существуют два набора инструментов, и оба бесплатны.

\subsection{Gowin EDA (официальная среда)}

Среду разработки \textbf{Gowin EDA Education} можно скачать с сайта
производителя (\url{https://www.gowinsemi.com}). Порядок работы:
\begin{enumerate}
  \item \textbf{File $\to$ New $\to$ FPGA Design Project}, указать имя проекта.
  \item Выбрать микросхему: серия \textbf{GW2AR}, корпус \textbf{QFN88},
        скорость \textbf{C8/I7}, устройство \textbf{GW2AR-18C}
        (полное имя \code{GW2AR-LV18QN88C8/I7}).
  \item \textbf{File $\to$ New $\to$ Verilog File} --- создать файл модуля
        и написать код.
  \item \textbf{File $\to$ New $\to$ Physical Constraints File} --- создать
        файл \code{.cst} с номерами выводов.
  \item Во вкладке \textbf{Process} запустить \textbf{Synthesize}, затем
        \textbf{Place \& Route}.
  \item Открыть \textbf{Programmer}, выбрать режим \textbf{SRAM Program}
        (для проверки) или \textbf{External Flash Mode} (чтобы схема
        сохранялась после выключения) и нажать \textbf{Program}.
\end{enumerate}

\subsection{Открытые инструменты}

Существует полностью открытый набор программ, работающий в Linux, Windows и macOS
(удобнее всего установить сборку \textbf{OSS CAD Suite}):
\begin{itemize}
  \item \textbf{Yosys} --- синтезатор;
  \item \textbf{nextpnr-himbaechel} --- размещение и трассировка;
  \item \textbf{Apicula} (\code{gowin\_pack}) --- упаковка битстрима для Gowin;
  \item \textbf{openFPGALoader} --- загрузка в плату;
  \item \textbf{Icarus Verilog} и \textbf{GTKWave} --- симуляция и
        просмотр временных диаграмм.
\end{itemize}

Все команды собраны в файл \code{Makefile} в папке с примерами:
\begin{vcode}[bash]
# 1. Синтез
yosys -p "read_verilog blink.v; synth_gowin -top blink -json blink.json"
# 2. Размещение и трассировка
nextpnr-himbaechel --json blink.json --write pnr.json \
    --device GW2AR-LV18QN88C8/I7 --vopt family=GW2A-18C \
    --vopt cst=blink.cst --freq 27
# 3. Битстрим
gowin_pack -d GW2A-18C -o blink.fs pnr.json
# 4. Загрузка: в SRAM ...
openFPGALoader -b tangnano20k blink.fs
# ... или во Flash
openFPGALoader -b tangnano20k -f blink.fs

# То же самое одной командой:
make PROJECT=02_blink TOP=blink load
\end{vcode}

\section{Файл ограничений .cst}
\label{sec:cst}

Код на Verilog ничего не знает о плате: в нём есть только имена портов
(\code{clk}, \code{led}, \code{btn}\ldots). Файл \code{.cst} (Physical
Constraints, «физические ограничения») сообщает трассировщику, к какой ножке
микросхемы подключить каждый порт и как эту ножку настроить.

\begin{figure}[H]
\centering
\begin{tikzpicture}[node distance=4mm and 20mm]
  \node[block, minimum width=30mm, minimum height=22mm, align=left, font=\ttfamily\scriptsize] (v) {module blink (\\\ \ input\ \ clk,\\\ \ output [5:0] led\\);};
  \node[oblock, right=of v, minimum width=36mm, minimum height=22mm, align=left, font=\ttfamily\scriptsize] (c) {IO\_LOC "clk" 4;\\IO\_LOC "led[0]" 15;\\\ldots\\IO\_LOC "led[5]" 20;};
  \node[draw=FpgaDark, fill=ChipBlack, text=white, minimum width=24mm, minimum height=22mm, right=of c, font=\sffamily\small, align=center] (chip) {ПЛИС\\\scriptsize выводы 4, 15..20};
  \draw[arr] (v) -- node[above, lbl] {имена} (c);
  \draw[arr] (c) -- node[above, lbl] {номера ножек} (chip);
  \node[lbl, below=1mm of v] {blink.v};
  \node[lbl, below=1mm of c] {blink.cst};
\end{tikzpicture}
\caption{Файл \code{.cst} связывает имена портов модуля с ножками микросхемы}
\end{figure}

В файле встречаются всего две команды:
\begin{description}
  \item[\code{IO\_LOC "порт" номер;}] --- подключить порт к выводу с указанным
    номером. Номера выводов Tang Nano 20K приведены в табл.~\ref{tab:pins}.
  \item[\code{IO\_PORT "порт" параметр=значение ...;}] --- задать
    электрические параметры вывода.
\end{description}

\begin{table}[H]
\centering
\caption{Основные параметры команды \code{IO\_PORT}}
\label{tab:ioport}
\small
\begin{tabularx}{\textwidth}{l l X}
\toprule
\textbf{Параметр} & \textbf{Значения} & \textbf{Смысл} \\
\midrule
\code{IO\_TYPE} & \code{LVCMOS33}, \code{LVCMOS18}, \ldots & стандарт логических уровней; на Tang Nano 20K почти везде \code{LVCMOS33} (3,3~В), см. главу~\ref{ch:signals} \\
\code{PULL\_MODE} & \code{UP}, \code{DOWN}, \code{NONE}, \code{KEEPER} & встроенный резистор: подтяжка к питанию, к земле, без подтяжки или «удержание» последнего уровня \\
\code{DRIVE} & \code{4}, \code{8}, \code{12}, \code{16}, \code{24} & максимальный выходной ток в миллиамперах (только для выходов) \\
\bottomrule
\end{tabularx}
\end{table}

Правила записи:
\begin{itemize}
  \item каждая команда заканчивается точкой с запятой;
  \item комментарии начинаются с \code{//}, как в Verilog;
  \item имена портов пишутся в кавычках и должны \emph{в точности} совпадать с
        именами в модуле (регистр букв тоже важен);
  \item у шины каждый бит описывается отдельной строкой: \code{"led[0]"},
        \code{"led[1]"}, \ldots
\end{itemize}

Зачем нужна подтяжка? Вход, к которому ничего не подключено (например, вывод
кнопки, пока она отпущена), «висит в воздухе» и ловит наводки: ПЛИС будет
читать то 0, то 1. Резистор \code{PULL\_MODE=DOWN} слабо притягивает такой вход
к земле, и пока кнопка отпущена, на нём гарантированно 0.

\begin{figure}[H]
\centering
\begin{circuitikz}[scale=0.85, transform shape]
  \draw (0,3.2) node[vcc] {$+3{,}3$ В} -- (0,2.8) to[push button, l=S1] (0,1.4) -- (2.4,1.4);
  \draw (1.3,1.4) to[R, l_=$R_\text{подтяжки}$, *-] (1.3,-0.2) node[ground] {};
  \draw (2.4,0.8) rectangle (4.6,2.0);
  \node[font=\sffamily\small] at (3.5,1.4) {ПЛИС};
  \node[font=\sffamily\scriptsize, anchor=south west] at (2.4,1.4) {88};
  \node[font=\sffamily\small, align=left, anchor=west] at (5.2,1.9) {кнопка отпущена: резистор тянет вход к 0};
  \node[font=\sffamily\small, align=left, anchor=west] at (5.2,1.0) {кнопка нажата: вход соединён с 3,3~В $\to$ 1};
\end{circuitikz}
\caption{\code{PULL\_MODE=DOWN}: резистор подтяжки включается внутри ПЛИС,
паять его не нужно}
\end{figure}

\subsection{Файлы .cst в примерах}

У каждого проекта из главы~\ref{ch:examples} есть свой файл ограничений
рядом с кодом, например \code{code/02\_blink/blink.cst}. В нём описаны ровно
те порты, которые есть в модуле, и каждая группа строк прокомментирована.
Вот полный файл для мигалки:

\vfile[c]{\codepath/02_blink/blink.cst}{code/02\_blink/blink.cst}

Кроме того, в \code{code/common/tang\_nano\_20k.cst} собрана сводная
распиновка всех выводов, которые используются в книге. Это справочник: не
подключайте его к проекту целиком, а копируйте из него нужные строки.

\begin{warnbox}[Частые ошибки]
Мы проверили поведение открытого трассировщика \code{nextpnr} на примерах:
\begin{itemize}
  \item \textbf{Порт модуля не описан в \code{.cst}} --- ошибка
        \code{Unconstrained IO}, прошивка не собирается. Так даже лучше: ошибка
        видна сразу.
  \item \textbf{В \code{.cst} описан порт, которого нет в модуле} --- строка
        просто пропускается с сообщением \code{Cell ... not found}, \emph{без
        ошибки}. Обычно это опечатка в имени: тогда настоящий порт остаётся
        неописанным, и срабатывает предыдущая ошибка. Увидев
        \code{Unconstrained IO}, первым делом ищите опечатку в \code{.cst}.
  \item \textbf{Неверный номер вывода} --- ошибка \code{Pin ... not found}.
\end{itemize}
\end{warnbox}

\subsection{Тактовая частота}

Трассировщику нужно знать, на какой частоте будет работать схема, чтобы
проверить, успевают ли сигналы пройти критический путь
(глава~\ref{ch:sync}). Для \code{nextpnr} частота задаётся ключом
\code{--freq 27} (в мегагерцах), и в журнале появляется строка вида
\begin{vcode}[text]
Info: Max frequency for clock 'clk': 256.21 MHz (PASS at 27.00 MHz)
\end{vcode}
В Gowin EDA частоту задают в отдельном файле временных ограничений \code{.sdc}:
\begin{vcode}[text]
create_clock -name clk -period 37.037 [get_ports {clk}]
\end{vcode}
Здесь 37,037~нс --- период сигнала 27~МГц.
